% ===================================================================
%  اللوحة نمرة ٨ — الحصاد.
%
%  Traduction, faite en 2026, en arabe egyptien ecrit, sur l'ido de
%  texto/io. Regles de la colonne : voir 10-lawha-01.tex. Deux
%  scenes, deux numerotations : les cles portent c1 et c2.
%
%  LA MOISSON DONNE A LA COLONNE SON PREMIER TROU, ET IL EST
%  INSTRUCTIF. Le fleau (40) — les deux batons articules dont on bat
%  le grain — n'a pas de nom egyptien, parce que l'Egypte ne bat pas
%  au fleau : elle bat au نورج, le traineau a rouleaux que les
%  bœufs trainent en rond sur l'aire. Nommer le fleau نورج aurait
%  change la CHOSE, ce que la colonne ne fait jamais ; on l'a donc
%  decrit — « عصي الدراس » — plutot que nomme. Sept tableaux ou
%  l'egyptien avait le mot d'avance ; celui-ci ne l'a pas, et il faut
%  le dire aussi.
%
%  POUR LE RESTE, LES OUTILS SONT LA :
%
%    منجل (25, 12)  la faucille
%    مقطف (17)      la hotte, le panier de dos
%    وابور (45)     la locomobile — le mot de toute machine a vapeur
%    قلاووظ (29)    la vis du pressoir
%    سنّارة (44)     la ligne du pecheur
%    فلوكة (49)     la barque du Nil
%    زمزمية (14)    la gourde
%
%  ET DEUX NOMS D'ETRES VIVANTS QUI VALENT LEUR LIGNE :
%
%    عين الجمل (57)  le noyer — « l'œil du chameau », pour la noix
%    أبو فصادة (19)  la bergeronnette
%
%  L'egyptien nomme volontiers par ressemblance et par surnom :
%  ابو فروة le marron au tableau 7, عين الجمل ici. Ce ne sont pas des
%  images litteraires, ce sont les noms du marche.
%
%  ECART DECLARE : le bloc c2-03-1 rend le renvoi (13) du gobelet,
%  que le fac-simile ido compose « 13) » sans parenthese ouvrante.
%  renvoji.py le passe, la raison est inscrite dans APARTA.
% ===================================================================

%%K t08-noto-1 noto
\VUnotes{143.2mm}{%
\textit{(*) علشان الكلمة دي مش دقيقة: هو ده حصاد كتّان، ولا قطف عنب، ولا جني
تفاح؟ يمكن يبقى أحسن نستعمل} \VUgras{mesono} \textit{المقترحة في «Mondo» الجزء
الرابعتاشر، صفحة ٢٤٢. بس لحدّ دلوقت الجذر الجديد ده لسه ما اتاخدش من المجمع،
وبيستنّى المناقشة على حسب القواعد المعتادة عند الإيديست.}%
}

%%K t08-apar-1 apar
\VUsaut{5.0mm}
\VUtitre{15.0pt}{60}{اللوحة نمرة ٨}
\VUsaut{3.0mm}
\VUcentre{13.2pt}{90}{\textit{المشهد الأول.}}
\VUsaut{3.0mm}
\VUpk{10.2pt}{40}{الحصاد (*).}
\VUsaut{4.0mm}
\VUpk{10.2pt}{40}{مناظر الريف.}
\VUsaut{3.0mm}

%%K t08-c1-01-1 p
1. --- مع \VUgras{الصيف}، شكل الريف اتغيّر مرة تانية. البذور اللي اتحطّت في
التلم نبتت؛ ونبات أخضر صغيّر طلع من الأرض وعمل سنابل. والحبّ اتكوّن، وبعدين
استوى، ومن دلوقت ما بقاش لازم غير الحصاد.

%%K t08-c1-02-1 p
2. --- \VUgras{ورد الغيط} فتّح. \VUgras{زهرة الذرة} \textsuperscript{(1)}،
و\VUgras{الخشخاش} \textsuperscript{(2)}، و\VUgras{كفّ الثعلب} \textsuperscript{(3)}
بيضيفوا للون غيطان القمح الواحد تباين \VUgras{تويجاتهم} الطازة المفرحة.

%%K t08-c1-03-1 p
3. --- وشجر الفاكهة اتغطّى بالورق؛ والفاكهة استوت. و\VUgras{شجرة الكرز}
\textsuperscript{(4)} الصغيرة مليانة \VUgras{كرز} \textsuperscript{(5)} حلو، والعيال
بيقطفوه وياكلوه بمتعة كبيرة.

%%K t08-c1-04-1 p
4. --- ومن دلوقت المواشي تقدر تقضّي النهار كله في الهوا الطلق.

%%K t08-c1-04-2 p
و\VUgras{الحملان} \textsuperscript{(6)} كبرت وبقت \VUgras{نعاج} \textsuperscript{(7)}
حلوة و\VUgras{خرفان} \textsuperscript{(8)} حلوة بصوف تخين. وهي بترعى بهدوء
\VUgras{حشيش} \VUgras{المرعى} \textsuperscript{(9)} المنقّط بـ\VUgras{الأقحوان}.

%%K t08-c1-04-3 p
وقريّب هيبدوا يجزّوا \VUgras{الحيوانات دي} علشان ياخدوا \VUgras{صوفها} اللي
بيتعمل منه قماش هدومنا.

%%K t08-tit-2 sub
\VUsaut{5.0mm}
\VUpk{10.2pt}{40}{الراعي.}
\VUsaut{3.4mm}

%%K t08-c1-05-1 p
5. --- \VUgras{الراعي} \textsuperscript{(10)} شكله مش واخد باله منهم أوي. هو مهتم بس
بالعاصفة اللي بتعدّي عند الأفق. و\VUgras{صبي الراعي} \textsuperscript{(13)}، اللي
بيقرّب \VUgras{الزمزمية} \textsuperscript{(14)} من بُقّه، شكله أقلّ اهتمامًا كمان.

%%K t08-c1-06-1 p
6. --- والحرّ خانق؛ علشان كده \VUgras{الكلب} \textsuperscript{(15)} جه، هو كمان،
يقطع عطشه؛ وهو بيلحس مية \VUgras{الترعة} \textsuperscript{(16)} الساقعة، خضّ
\VUgras{ضفدعة} \textsuperscript{(17)} مسكينة كانت مستخبّية في حشيش الشطّ. ودي نطّت
في المية الصافية اللي بيزهّر فيها \VUgras{زنبق المية} \textsuperscript{(18)}.

%%K t08-c1-07-1 p
7. --- و\VUgras{أبو فصادة} \textsuperscript{(19)} بيستحمّى جنب \VUgras{العين}
\textsuperscript{(20)} اللي بتنطّ بين حجرتين، وبيستخبّى تحت \VUgras{الشجيرات
الشوكية} \textsuperscript{(21)} و\VUgras{شجر الزعرور} \textsuperscript{(22)}.

%%K t08-tit-3 sub
\VUsaut{5.0mm}
\VUpk{10.2pt}{40}{الحصاد.}
\VUsaut{3.4mm}

%%K t08-c1-08-1 p
8. --- بالنسبة لعيال الريف، حصاد القمح وقت مسلّي أوي. وهما \VUgras{بيشقلبوا}
\textsuperscript{(24)} بفرحة على \VUgras{أكوام القش} \textsuperscript{(23)}، وبيساعدوا
\VUgras{الحصّادين} \textsuperscript{(26)}. ودول بيحصدوا \VUgras{غيط القمح}
\textsuperscript{(27)}، و\VUgras{مناجلهم} \textsuperscript{(25)} مسنونة كويس
بـ\VUgras{حجر السنّ} \textsuperscript{(30)}؛ والعيال بيلمّوا القمح ويربطوه
\VUgras{حزم} \textsuperscript{(31)} أو \VUgras{ربطات} \textsuperscript{(32)}.

%%K t08-c1-09-1 p
9. --- وساعات بيسمحوا لأكبر واحد فيهم إنه يقطع بـ\VUgras{منجل} \textsuperscript{(12)}
\VUgras{السيقان الدهبية} \textsuperscript{(28)} اللي شايلة \VUgras{السنابل}
\textsuperscript{(29)} التقيلة المليانة حبّ؛ والصغيّرين، وهما بيراقبوا
\VUgras{المواشي} \textsuperscript{(33)} اللي بترعى في \VUgras{المرعى}
\textsuperscript{(34)} تحت ضلّ \VUgras{شجرة الزيزفون} \textsuperscript{(35)} الكبيرة،
بيمسكوا بـ\VUgras{شبكتهم} \textsuperscript{(36)} \VUgras{الفراشات} الحلوة.

%%K t08-c1-09-2 p
وفي منهم اللي بيطلع على \VUgras{العربية} \textsuperscript{(37)} علشان يرتّب الحزم
اللي بيناولوها له بـ\VUgras{المذراة} \textsuperscript{(38)}.

%%K t08-c1-10-1 p
10. --- وعلى \VUgras{السهل} \textsuperscript{(39)} كله، اللي \VUgras{القبّرات}
\textsuperscript{(41)} بتطير فيه، الحركة مالياه. الكل مشغول بالحصاد. بس، وفي الوقت
اللي الفقرا فيه لسه بيستعملوا العدّة القديمة، \VUgras{المناجل}، و\VUgras{عصي
الدراس} \textsuperscript{(40)}، الملّاك الأغنيا بيحصدوا قمحهم أو \VUgras{شيلمهم}
\textsuperscript{(43)} بـ\VUgras{ماكينة الحصاد} \textsuperscript{(42)}. وبعد كده، من
غير ما يودّوا الحزم للجرن، بيعملوا \VUgras{أكوام قمح} \textsuperscript{(44)} كبيرة.

%%K t08-c1-11-1 p
11. --- وساعتها بيجيبوا \VUgras{ماكينة الدراس} \textsuperscript{(47)} اللي
\VUgras{الوابور} \textsuperscript{(45)} بيحرّكها بـ\VUgras{سير نقل} \textsuperscript{(46)}؛
وكده الحبّ بينفصل عن \VUgras{القش} من غير ما يسيب الغيط اللي اتزرع فيه.

%%K t08-tit-4 sub
\VUsaut{5.0mm}
\VUpk{10.2pt}{40}{العاصفة.}
\VUsaut{3.4mm}

%%K t08-c1-12-1 p
12. --- وكتير بتحصل \VUgras{عواصف} \textsuperscript{(53)} وقت الحصاد. الأول بيتسمع
هدير \VUgras{الرعد} من بعيد؛ وبعدين \VUgras{ريح غربية} \textsuperscript{(54)} قوية
بتهبّ وتحني وهي بتلهث أتخن شجر \VUgras{الغابة} \textsuperscript{(49)}.
و\VUgras{السحب السودا} \textsuperscript{(56)} بتهجم على السما؛ وبيخترقها زجزاج
\VUgras{البرق} \textsuperscript{(55)} اللي بيعمي. و\VUgras{المطر} وأحيانًا
\VUgras{البَرَد} بينزلوا ويفرشوا المحصول اللي كان جاهز للحصاد.

%%K t08-c1-13-1 p
13. --- والبرق كتير بيولّع المباني. وكده السنة اللي فاتت سقف \VUgras{طاحونة
الهوا} \textsuperscript{(50)} بقى رماد واتكسر اتنين من \VUgras{أجنحتها}
\textsuperscript{(51)}.

%%K t08-c1-14-1 p
14. --- وبعد كام ساعة، الهدوء بيرجع. و\VUgras{قوس قزح} \textsuperscript{(52)} لمّاع
بيبان عند الأفق، والطبيعة بترجع لحياتها اللي اتقطعت كام لحظة. وأهل الريف، اللي
كانوا لجأوا لـ\VUgras{العشش} \textsuperscript{(48)}، بيرجعوا للشغل ويجتهدوا بشجاعة
يصلّحوا الخسارة اللي حصلت لهم.

%%K t08-tit-5 sub
\VUsaut{6.0mm}
\VUcentre{13.2pt}{90}{\textit{المشهد التاني.}}
\VUsaut{3.6mm}

%%K t08-tit-6 sub
\VUpk{10.2pt}{40}{الصيد. --- قطف العنب.}
\VUsaut{1.4mm}
\VUpk{10.2pt}{40}{صيد السمك.}
\VUsaut{4.0mm}

%%K t08-c2-01-1 p
1. --- في آخر \VUgras{أغسطس} تقريبًا أو نصّ \VUgras{سبتمبر}، على حسب المناطق،
بيبتدي موسم الصيد. وأول شعاع شمس بالكاد بيخلّي \VUgras{السحالي} \textsuperscript{(2)}
تخرج من جحورها، لمّا \VUgras{الصيّاد} \textsuperscript{(5)} بيبتدي طريقه مع
\VUgras{كلابه} \textsuperscript{(4)}.

%%K t08-c2-02-1 p
2. --- لبس \VUgras{بدلة الصيد} \textsuperscript{(9)} المتينة من قماش تخين، ولفّ على
رجليه \VUgras{لفّات} جلد، اللي بيها هيقدر يمشي في الحشيش المبلّل اللي بتستخبّى
فيه \VUgras{العلاجيم} \textsuperscript{(1)}؛ وخد \VUgras{بندقيته} \textsuperscript{(6)}
و\VUgras{شنطة الصيد} \textsuperscript{(10)} بتاعته، وأهو ماشي.

%%K t08-c2-03-1 p
3. --- وحماسه في الجري ورا الطرايد وصّله لوسط حقل عنب، اللي فيه قطف العنب
شغّال بحيوية. وعيال بتاع العنب، اللي عادةً بيراقبوا المعزات على الطريق،
النهارده سايبينهم يرعوا على مزاجهم؛ وبعد ما ربطوا في وتد \VUgras{تيس}
\textsuperscript{(3)} عجوز، اللي أكيد مش هيفوّت فرصة حرّيته ويهرب، خدوا لنفسهم
هما كمان نصيب من المتعة. واحد فيهم خد عنقود عنب من \VUgras{سبت القطف}
\textsuperscript{(12)}، وفضل يسلّي نفسه بإنه يصبّ \VUgras{العصير} \textsuperscript{(14)}
في \VUgras{كباية} \textsuperscript{(13)} علشان يعمل عصير عنب (نبيت ما اتخمّرش).
والتاني حطّ على راسه \VUgras{تاج ورق} \textsuperscript{(15)} كان حالًا ضفره.

%%K t08-c2-03-1 p suite
وجنبهم، \VUgras{عصافير} \textsuperscript{(16)} قليلة الأدب بتيجي لحدّ قدام المعصرة
علشان تسرق العنب.

%%K t08-c2-04-1 p
4. --- ولمّا العيال بتلعب، الرجّالة بتشتغل. \VUgras{قاطفين العنب}
\textsuperscript{(18)} بيودّوا العنب للقبو في \VUgras{مقاطف} \textsuperscript{(17)}؛
و\VUgras{بتاع البراميل} \textsuperscript{(19)} بيصلّح \VUgras{البراميل}
\textsuperscript{(20)}، ويتأكّد إن \VUgras{ألواحها} \textsuperscript{(22)}
و\VUgras{قيعانها} \textsuperscript{(21)} كويسة، ويغيّر \VUgras{الأطواق}
\textsuperscript{(23)} المكسورة. وهو كمان اشتغل \VUgras{الأحواض الكبيرة}
\textsuperscript{(25)} اللي بيصبّوا فيها \VUgras{العنب} \textsuperscript{(26)} علشان
يتخمّر.

%%K t08-c2-05-1 p
5. --- ولمّا التخمير بيخلص، بيسحبوا النبيت ويحطّوا البقايا على \VUgras{المعصرة}
\textsuperscript{(28)}، وبالراحة، وهما بيلزقوا بـ\VUgras{القلاووظ} \textsuperscript{(29)}،
بيعصروا العصير اللي بيسيل في \VUgras{حوض صغير} \textsuperscript{(32)}، وبيطلّعوا
كمان \VUgras{نبيت} \textsuperscript{(31)}.

%%K t08-tit-8 sub
\VUsaut{6.0mm}
\VUpk{10.2pt}{40}{البيرة وعصير التفاح.}
\VUsaut{3.4mm}

%%K t08-c2-06-1 p
6. --- وعلى حيطة \VUgras{القبو} \textsuperscript{(27)} متسلّق فرع \VUgras{جنجل}
\textsuperscript{(30)}. وهنا هو نبات زينة بس، إنما في بلاد فيها العنب نادر أوي،
بيزرعوا الجنجل علشان يعملوا \VUgras{البيرة}.

%%K t08-c2-07-1 p
7. --- و\VUgras{شجرة التفاح} \textsuperscript{(33)} الصغيرة مليانة تفاح، اللي هيطلّع
لمّا يستوي \VUgras{عصير تفاح} ممتاز. وفي \VUgras{قفصها} \textsuperscript{(34)}،
\VUgras{الكناريا} \textsuperscript{(35)} و\VUgras{الحسّون} \textsuperscript{(36)}
بيبصّوا بعين الغيرة للعصافير اللي بتلعب بحرّية.

%%K t08-c2-08-1 p
8. --- ولحدّ آخر المدّ البصر، على منحدر التلال، متصفّفة \VUgras{عيدان العنب}
\textsuperscript{(40)} اللي شايلة \VUgras{جذوع العنب} \textsuperscript{(39)} الحلوة أوي،
المزيّنة بـ\VUgras{عناقيد} تقيلة.

%%K t08-c2-08-2 p
وطول السنة، \VUgras{بتاع العنب} \textsuperscript{(42)} اشتغل علشان يعتني
بـ\VUgras{حقل العنب} \textsuperscript{(38)} بتاعه، ودلوقت

%%K t08-c2-08-2 p suite
بياخد أجر تعبه محصول حلو. و\VUgras{قاطفات العنب} \textsuperscript{(41)} بيملوا
الأحواض المحطوطة على العربية اللي بتجرّها \VUgras{البغال} \textsuperscript{(43)}،
والكل مبسوط.

%%K t08-tit-9 sub
\VUsaut{5.0mm}
\VUpk{10.2pt}{40}{صيد السمك.}
\VUsaut{3.4mm}

%%K t08-c2-09-1 p
9. --- ونفس الكلام على \VUgras{الصيّادين بالسنّارة} \textsuperscript{(44)}، اللي من
الصبح جايين قاعدين على شطّ المية بـ\VUgras{سنانيرهم} \textsuperscript{(45)}.

%%K t08-c2-09-2 p
و\VUgras{السمك} \textsuperscript{(47)} بيخطف \VUgras{الشُصّ} أول ما ينزّلوا
\VUgras{خيطهم} \textsuperscript{(46)} في المية، وبالكاد بيلحقوا يغيّروا
\VUgras{الطُّعم}، لحدّ ما ضحية جديدة تنطّ في آخر الخيط.

%%K t08-c2-09-3 p
ومع كده \VUgras{الصيّاد بالشبكة} \textsuperscript{(48)} اللي بيرمي \VUgras{الشبكة}
من \VUgras{فلوكته} \textsuperscript{(49)} لوسط \VUgras{النهر} \textsuperscript{(50)}
بيمسك سمك أكتر.

%%K t08-c2-10-1 p
10. --- والخريف هو فصل النزهات الحلوة: على الرجلين، أو \VUgras{بالحصان}
\textsuperscript{(52)}، أو \VUgras{بالعجلة} \textsuperscript{(53)}، أو
\VUgras{بالعربية} \textsuperscript{(54)}. و\VUgras{الطرق} \textsuperscript{(51)} نضيفة
والناس بتستغلّ آخر أيام الجو الحلو، علشان أهو \VUgras{السنونو} \textsuperscript{(56)}
بيتجمّع على السقوف علشان يطير بكل جناحه لبلاد أدفا. وورق \VUgras{شجرة عين
الجمل} \textsuperscript{(57)} العجوزة بيصفرّ، و\VUgras{عين الجمل} بيقع، والشجر
العالي المتجمّع زي \VUgras{الباقة} \textsuperscript{(58)} على \VUgras{المرتفع}
\textsuperscript{(59)} قريّب هيفضى من ورقه.

%%K t08-c2-11-1 p
11. --- \VUgras{الشمس} \textsuperscript{(60)} بتطلع متأخّر، والنهار بيقصر، وبرد
قارص شوية بيبتدي يتحسّ \VUgras{الصبح}، وأهو أسراب \VUgras{الشنقب}
\textsuperscript{(61)} بتسيب مناخنا اللي ما بقاش بيرحّب بيها.

\VUsaut{6.0mm}
\VUfilet{22mm}
